% !Mode:: "TeX:UTF-8"
% !TEX program  = xelatex

\section{引言}

旅行商问题（Traveling Salesman Problem, TSP）和容量约束车辆路径问题（Capacitated Vehicle Routing Problem, CVRP）
是组合优化中两类经典问题。近年来，神经组合优化（Neural Combinatorial Optimization, NCO）提出了多种方法\cite{bello2016nco,kool2019attention}，
但经验表明，单一方法很难在不同实例和不同分布上始终保持最优。

首先要关注的问题是：实例级的方法选择是否真的必要？
预分析统计表明，无论在合成分布还是在 LIB 基准数据集上，都不存在一种方法能够在同一测试集上稳定支配其他方法：
例如在 61 个 TSPLIB 实例上，LEHD 虽然胜出最多（70.5\%），但仍有近 30\% 的实例被其他方法超越；方法偏好也不能仅靠规模 $n$ 这类单一特征解释（详细统计见附录 A）。
这一实验说明，不同方法在不同实例上具有显著互补性，实例级选择存在切实的改进空间。

尽管 Gao et al.\cite{gao2025nss} 已经初步验证了方法选择的有效性，
但该工作采用全信息监督学习，且对每个问题类型分别训练独立模型。
本文聚焦 TSP 和 CVRP 的初始化阶段，研究如何依据实例特征从固定候选池中选择一种初始化方法，使初始解质量最优。
本文将这一任务建模为单步上下文老虎机问题，以更贴近实际应用中在线的逐实例做出选择、逐步收集反馈的场景。
围绕这一建模，需要解决两个关键问题：
（1）如何在 TSP 与 CVRP 两类问题之间构造共享的动作空间并统一表示实例特征？
（2）如何设计奖励（reward）函数，使得在两种问题上都有效？

本文提出基于 Neural-LinUCB 的统一算子选择模型，在单一模型中同时处理 TSP 和 CVRP 的初始化方法选择。
本文的主要贡献如下：
\begin{itemize}
  \item 将神经组合优化中的初始化方法选择形式化为单步上下文老虎机问题，并给出统一的动作空间设计和 reward 定义。
  \item 提出基于 Neural-LinUCB 框架的统一选择模型，融合层次化图编码、手工统计特征与可学习方法嵌入构造动作条件表征，并通过按问题拆分的浅层 UCB 头平衡利用与探索。该模型以单一架构同时服务 TSP 和 CVRP。
  \item 在合成混合数据集上验证 Neural-LinUCB 取得优于固定最佳单方法的 mean length，同时通过分析发现当前 reward 设计与最终长度目标之间的偏差以及分布外泛化的不足。
\end{itemize}

本文其余部分组织如下：第 \ref{sec:related} 节介绍相关工作；第 \ref{sec:method} 节给出问题定义与方法设计；第 \ref{sec:experiment} 节说明实验设置；第 \ref{sec:results} 节报告结果与分析；第 \ref{sec:conclusion} 节总结全文。
